\documentclass[11pt,a4paper]{article}
\usepackage[margin=2.2cm]{geometry}
\usepackage{amsmath,amssymb}
\usepackage[T1]{fontenc}\usepackage{mathptmx}
\pdfminorversion=4
\newcommand{\rev}[1]{\par\medskip\noindent\textbf{Reviewer:} \emph{#1}\par\smallskip\noindent\textbf{Response:} }
\begin{document}
\begin{center}
{\Large\bfseries Response to the Reviewer (Third Round)}\\[4pt]
Flatness and Passivity-Based Control with Adaptive Neural Compensation of a Direct-Drive PMSG Wind Energy Conversion System\\[2pt]
S. Regani, M. Messadi, K. Kemih
\end{center}

We thank the reviewer for the careful re-evaluation. The remaining clarifications have been made as
follows; the paper keeps its six-page length (the abstract and one paragraph of Section~VI-C were
condensed).

\rev{1. The local nature of the stability guarantees should be emphasized, particularly regarding actuator saturation and model uncertainties.}
A paragraph was added after Proposition~3 (Section~V-C). It states that Propositions~2--3 are local
guarantees, which assume:
\begin{itemize}
\item an exact model of the electrical part;
\item no saturation;
\item trajectories that remain in the operating region.
\end{itemize}
It also states that the behavior with saturations (start-up, grid dip, $\Gamma=20$, initial speed
errors) and with parameter mismatch is assessed only numerically, in Section~VI. The conclusion lists
an analysis including saturations as future work.

\rev{2. The justification for neglecting the energy-reference derivative should explicitly distinguish steady-state and transient operating conditions.}
Section~V-A now states explicitly that the bound of 1.4~V holds only in steady operation, where
$|\dot y_2^*|\le2.7$~kW. In the transients, $|\dot y_2^*|$ reaches 29~kW at start-up and 52~kW during
the grid dip. The same $L_1$ bound then gives conservative values of 15 and 27~V. The simulated
deviations, which include this effect, are 10.4 and 3.6~V.

\rev{3. The interpretation of IEEE 519 harmonic limits should clearly distinguish converter-terminal measurements from compliance at the point of common coupling.}
Section~VI-D now states that the TDD values are computed at the converter terminals of a single unit
and do not by themselves establish compliance at the PCC. Compliance at the PCC also depends on the
short-circuit ratio, the transformer, the collector network and the other units. The text also keeps
the switching interharmonics, and the resulting need for an $LCL$ filter, separate from the harmonic
indices ("Separately, \dots"; "independently of the harmonic indices, \dots").

\rev{4. Reference [15] should be updated to reflect its current publication status.}
At the time of this revision, [15] is still under review at \emph{Mathematics} and has no DOI. It is
cited as such. It will be replaced by the full citation with its DOI as soon as it is accepted, and
in any case before the camera-ready version. No stability result of the paper depends on it.

\rev{Minor remark 4: balanced presentation of the benefits.}
The abstract and the conclusion now state both sides:
\begin{itemize}
\item the proposed control improves the MPPT tracking (3.8 times smaller IAE) and the DC-link
behavior during the dip (3.6 vs 11.6~V);
\item PI gives smaller RMS current errors (0.84 vs 3.07~A for $i_{sq}$) and, with feedforward, a smaller
start-up DC-link deviation (3.9 vs 10.4~V).
\end{itemize}
\end{document}
